%
% synopsis.tex — Tête de veau ravigote · Éric Mugnier
% Compilation : latexmk -pdf synopsis.tex
%
\pdfminorversion=7
\documentclass[12pt,a4paper]{article}

\usepackage[utf8]{inputenc}
\usepackage[T1]{fontenc}
\usepackage{ebgaramond}
\usepackage[french]{babel}
\usepackage{microtype}
\usepackage[left=30mm,right=25mm,top=28mm,bottom=25mm]{geometry}

\setlength{\parindent}{0pt}
\setlength{\parskip}{7pt}

\begin{document}
\pagestyle{empty}

\begin{center}
{\large\textsc{éric mugnier}}

\bigskip

{\Large\textit{Tête de veau ravigote}}

\medskip

{\small roman}
\end{center}

\bigskip\bigskip

\begin{quote}
Quand le père Vidal disparaît dans des circonstances troubles\\
et que des colis macabres commencent à circuler dans la ville,\\
le commandant Beauvais se lance dans une enquête labyrinthique\\
qui le mènera des caves de l'Église à l'ambassade de Namibie.

\medskip
Entouré de Titus Beaugendre, son fidèle équipier,\\
et d'une galerie de personnages hauts en couleur,\\
Beauvais arpente un monde contemporain\\
où l'absurde le dispute à l'horreur.
\end{quote}
\medskip

Un ovni, plutôt qu'un polar ; une logorrhée vengeresse et jubilatoire, écrite en
\textit{spontaneous prose}, où l'humour potache s'achoppe à des scènes \textit{gore}
à la limite du soutenable ; un roman-fleuve traversé d'amples digressions en forme
de cabinet de curiosités : sur les chats, les dinosaures, l'Église, les pavillons
de banlieue, l'unité 731, Pannonica de Koenigswarter, la 'Ndrangheta.

L'\textbf{Horreur} (prononcé à la Marlon Brando dans Apocalypse Now) se détache
au fil des pages comme le thème central du livre.

Toute forme de morale est étrangère aux protagonistes — au héros-narrateur en premier.
Entre grotesque et tragique, les digressions dessinent par intermittence l'image d'une
humanité sans Dieu, désespérée. Le ton n'est jamais didactique, jamais emphatique,
toujours factuel, genre \og Alain Decaux raconte \fg{}. Narrer L'Horreur sur ce ton
d'entomologiste provoque spontanément une sensation d'ironie dans l'esprit du lecteur,
un mécanisme inné de défense de son intégrité mentale. Bien plus efficace qu'un bouquin
qui dirait benoîtement : \og Regardez ça, c'est mal ! Ouh comme c'est vilain ! \fg{}.
Un bruit de fond, un acouphène, un filtre optique polarisé, un voile de gaze sur la peau.
La tentaculaire progression de la Camorra dans le Naples de l'\textit{Amica Geniale}
d'Elena Ferrante, jamais explicite, trop prégnante.

\end{document}
